\section{Training Agentic Visual Generation}
\label{sec:training}

Training an \AVG{} system is not equivalent to improving a visual generator in isolation. The generator determines which artifacts are reachable. The agentic loop determines how goals are decomposed, tools are selected, evidence is inspected, failures are diagnosed, results are revised, and execution is stopped. The central training problem is therefore to improve the policy that maps an evolving visual state, execution history, and user interaction to the next action. Fine-tuning a diffusion model, code model, or editor changes the executor's capabilities and the set of reachable artifacts. It constitutes agentic training only when the update also changes state-dependent decisions or the behavior of the loop.

This distinction separates three practices that recent studies often place under the same label. \textbf{Direct policy training} updates a controller, planner, router, critic, or integrated multimodal policy through trajectories, preferences, or reinforcement learning~\cite{avg260328767,avg260327742,jiang2026genagent,avg260302681,avg260517969,avg260308059,avg260705465}. \textbf{Agentic data curation} uses agents to construct data for a conventional visual model~\cite{avg260320644,avg260603168,avg260631537}. \textbf{Persistent non-parametric adaptation} freezes model parameters but updates persistent memory, skills, or routing knowledge from accumulated experience~\cite{avg260628971,avg260202051,jiang2026octot2i,avg260701709,avg260522208}. All three contribute to capability acquisition, but they support different claims. Agentic data curation does not show that the curator learned to act. Likewise, retained experience does not make a training-free system parametrically trained.

The chapter therefore follows the causal path by which experience changes future behavior. Section~\ref{subsec:training-targets} identifies the trainable objects and establishes the boundary of agentic training. Section~\ref{subsec:trajectory-supervision} examines the construction of visual trajectories and the supervision they preserve. Section~\ref{subsec:policy-initialization} studies supervised and preference-based policy initialization, and Section~\ref{subsec:multimodal-rl} turns to reinforcement learning from multimodal and executable feedback. Section~\ref{subsec:experience-distillation} distinguishes parametric experience distillation from non-parametric memory and skill updates. Section~\ref{subsec:training-transfer} then considers joint optimization, transfer, forgetting, provenance, and the evidence required to substantiate cross-task improvement.

\begin{keytakeaway}[Training Agentic Visual Generation at a Glance]
\begin{keypoints}
  \item \textbf{Training targets and boundaries:} identify which component or persistent knowledge store is changed and whether the update changes observation-conditioned decisions.

  \item \textbf{Visual trajectory data and supervision:} preserve the task, pre-action state, action, resulting observation, critique or reward, and updated plan or stopping decision.

  \item \textbf{Supervised and preference-based policy initialization:} establish executable behavior and learn relative trade-offs before interaction-based optimization.

  \item \textbf{Reinforcement learning:} connect decisions to multimodal and executable consequences, with rewards matched to the component being trained.

  \item \textbf{Experience reuse, skill abstraction, and distillation:} distinguish episodic retrieval, reusable skills, non-parametric memory updates, and parametric distillation.

  \item \textbf{Joint optimization, transfer, and evidence:} isolate the updated component and test transfer, forgetting, and persistence under controlled budgets and independent evaluation.
\end{keypoints}
\end{keytakeaway}

\subsection{Training Targets and Boundaries}
\label{subsec:training-targets}

The first question is not which optimization algorithm is used, but which component or persistent knowledge store is changed. In AVG, updates may target the generator, controller, router, critic, memory, or an integrated policy. Each object governs a different part of system behavior and therefore requires distinct data, objectives, and evidence. Generator updates can improve execution quality without changing tool selection. \textbf{Controller updates}, by contrast, can reduce redundant actions without changing the generator. Verifier updates reshape the feedback or reward landscape, while memory and skill updates can alter later behavior without modifying model parameters. Table~\ref{tab:training-targets} summarizes these causal roles.

\begin{table}[t]
  \caption{Trainable and updatable objects in \AVG{}.}
  \label{tab:training-targets}
  \small
  \setlength{\tabcolsep}{5pt}
  \renewcommand{\arraystretch}{1.08}
  \renewcommand{\tabularxcolumn}[1]{m{#1}}
  \begin{tabularx}{\textwidth}{@{}
    >{\raggedright\arraybackslash\bfseries}m{0.23\textwidth}
    >{\raggedright\arraybackslash}m{0.42\textwidth}
    >{\raggedright\arraybackslash}X@{}}
    \toprule
    Learned or updated object & \textbf{Behavior changed} & \textbf{Form of persistence} \\
    \midrule
    Generator or executor & Rendering, editing, code generation, and specialized tool execution & Generator, editor, code-model, or tool parameters \\
    \midrule
    Controller or planner & Goal decomposition; state interpretation; action ordering; revision, clarification, and stopping & Policy-model parameters or a structured planning policy \\
    \midrule
    Router or coordinator & Selection and assignment of generators, retrieval sources, tools, or specialized agents & Router parameters, rules, or routing memory \\
    \midrule
    Critic, verifier, or reward model & Quality assessment; error localization; preference prediction; feedback generation & Critic or evaluator parameters and calibration state \\
    \midrule
    Memory or skill system & Retrieval of prior attempts; workflow reuse; routing adaptation; procedural abstraction & Episodic memory, skill library, workflow template, or routing memory \\
    \midrule
    Integrated native policy & Joint understanding, deliberation, planning, and creation & Shared multimodal policy parameters \\
    \bottomrule
  \end{tabularx}
\end{table}
\FloatBarrier

Updates to the controller, or to the decision-making component of an integrated policy, provide the most direct evidence of agentic training under the survey's operational definition. GenAgent first establishes tool use through supervised fine-tuning, then optimizes multiround decisions with outcome and process rewards~\cite{jiang2026genagent}. TIR-Agent learns when to inspect, which restoration tool to apply, and when to stop~\cite{avg260327742}. Generation Navigator formulates text-to-image generation as state-conditioned action selection~\cite{avg260517969}. CanvasAgent learns executable reasoning-action trajectories and adapts tool choices after observing intermediate canvas states~\cite{avg260705465}. These systems train the mapping from observed state to subsequent action rather than the visual backbone alone.

\textbf{Generator training} has a neighboring but distinct causal role. A stronger executor enlarges the set of reachable artifacts, but an unchanged controller may still choose poor tools or compose them badly. OPERA addresses this mismatch by optimizing a restoration planner and co-training tools under the plans in which they are composed~\cite{avg260522104}. This coupling provides stronger evidence of agentic training than independent tool fine-tuning. The Manim study similarly separates code-generation improvements from visually grounded rewards on rendered animations~\cite{avg260418364}. Executor training can contribute to the training of an AVG system when it is coupled to the agentic loop, as in plan-conditioned or jointly optimized execution. However, executor improvement alone is not evidence that the agentic policy has learned. Stronger evidence requires the update to change observation-conditioned decisions or to be jointly optimized with the controller under loop-level feedback. Ordinary visual-model pretraining therefore remains background.

\textbf{Agentic data curation} follows the same boundary. ScaleEdit-12M, DataEvolver, and JAVEDIT use agentic workflows to construct, filter, or verify visual training data~\cite{avg260320644,avg260603168,avg260631537}. These systems can improve supervision for a downstream generator, but they do not show that the data-producing agent acquires a reusable policy. Section~\ref{subsec:trajectory-supervision} therefore treats them as data sources and as cases of provenance and filtering risk, not as direct evidence of controller learning.

\textbf{Persistent non-parametric adaptation} must be separated from parametric training and within-episode adaptation. OctoT2I updates routing knowledge without supervised fine-tuning~\cite{jiang2026octot2i}. COMFYCLAW converts trajectories and feedback into reusable skills~\cite{avg260701709}, while restoration systems retain high-value experience for later tasks~\cite{avg260628971,avg260202051,avg260522208}. Such persistent updates support stronger autonomy claims only when they change behavior on held-out future tasks. A longer context, extra reflection, or disposable search procedure demonstrates local adaptation rather than cross-task learning. CRAFT, VISTA, and Qwen-Image-Agent are useful boundary cases because they emphasize training-free or test-time operation~\cite{kovalev2025craft,long2026vista,avg260626907}. Section~\ref{subsec:experience-distillation} specifies the evidence needed to distinguish episodic reuse, persistent adaptation, and continual self-improvement.

\subsection{Visual Trajectory Data and Supervision}
\label{subsec:trajectory-supervision}

Direct training of agentic policies benefits from supervision that distinguishes decisions and their consequences rather than specifying desirable final artifacts alone. A prompt--image pair defines an acceptable endpoint but does not reveal the evidence behind tool selection, revision, or stopping. A visual trajectory can instead preserve the task, pre-action state, selected action and arguments, resulting observation, critique or reward, and updated plan or stopping decision. Preserving these dependencies makes it possible to learn how behavior should change as the artifact evolves, rather than merely imitate successful output sequences.

\textbf{Successful executable trajectories} provide a practical starting point because they preserve both action syntax and observable consequences. CanvasAgent learns multiround orchestration from fully annotated CanvasCraft trajectories before reinforcement learning~\cite{avg260705465}. TOOLCAD similarly grounds actions in a CAD environment with executable intermediate operations~\cite{gong2026toolcad}. Success-only traces, however, introduce selection bias: they show what worked but provide limited evidence about how to recognize deterioration, recover from failed calls, or distinguish necessary from redundant actions. \textbf{Failure--diagnosis--repair trajectories} are therefore especially useful when the target behavior includes reflection, correction, rollback, or adaptive stopping.

\textbf{Pairwise trajectory construction} provides complementary supervision by exposing relative decision quality. GenAgent compares final images and consecutive rounds while filtering revisions that reduce quality~\cite{jiang2026genagent}. GenEvolve contrasts tool-orchestrated trajectories for the same request and structures their differences as reusable experience~\cite{chen2026genevolve}. Clarify Before Executing similarly selects high-scoring trajectories and derives preference pairs from them~\cite{avg260716352}. Such data make the consequences of alternative decisions more explicit, although their usefulness still depends on the reliability and independence of the evaluator used to construct the comparison.

Long visual trajectories also impose a systems constraint. Each rollout may contain large intermediate images, repeated generator calls, executable states, and multimodal critiques. InterleaveThinker therefore decomposes expensive interleaved trajectories into single-step planner and critic data with step-wise rewards~\cite{avg260613679}. VisionCreator instead uses a simulated environment to reduce repeated interaction with external tools~\cite{avg260302681}. Both strategies improve scalability, but they discard different information: single-step optimization weakens full-trajectory credit assignment, whereas simulation introduces environment mismatch. Reliable comparison therefore benefits from reporting trajectory length, generator calls, simulator assumptions, filtering rates, and the retained distribution of failures rather than only prompt or image counts.

Agentic data curation introduces a related but distinct source of supervision. ScaleEdit-12M, DataEvolver, and JAVEDIT use agents and task-aware verifiers to construct, filter, or validate visual training data~\cite{avg260320644,avg260603168,avg260631537}. Such pipelines can substantially expand downstream supervision, but retrieved or generated assets and model-based judgments may also introduce provenance, licensing, evaluator-circularity, and leakage risks. Evidence is therefore stronger when data sources, model and tool versions, filtering and deduplication procedures, and overlap between training and evaluation pipelines are made explicit. Otherwise, additional scale may amplify shared evaluator bias or benchmark-specific patterns rather than improve transferable decisions.

\subsection{Supervised and Preference-Based Policy Initialization}
\label{subsec:policy-initialization}

\textbf{Supervised fine-tuning (SFT)} primarily establishes an executable initial policy. It can teach valid action syntax, plan structure, state-referenced tool arguments, observation-conditioned reasoning, and plausible stopping behavior. GenAgent uses tool-use and reflection supervision before reinforcement learning~\cite{jiang2026genagent}. CanvasAgent first learns executable reasoning--action trajectories~\cite{avg260705465}, while TIR-Agent uses SFT before optimizing exploration and efficiency~\cite{avg260327742}. Across these systems, supervised initialization stabilizes the interface between reasoning and visual execution, but does not by itself guarantee effective decisions outside the demonstration distribution.

The composition and ordering of supervised data can matter as much as its volume. VisionCreator argues that generation-only SFT may erode broader multimodal ability, whereas indiscriminate mixture training may underdevelop creation expertise. Its progressive specialization curriculum preserves general capability while preparing the policy for later reinforcement learning~\cite{avg260302681}. VisionCreator-R1 provides complementary supervision for planning and reflection before multitask reinforcement learning~\cite{avg260308812}. IEA combines expert-distilled SFT, group relative policy optimization, and synthetic fine-tuning~\cite{avg260608016}. Together, these systems show that supervised alignment may be organized as a staged curriculum and, in some cases, revisited after reinforcement learning rather than serving only as a one-time initialization stage.

Decomposition provides another way to make visual policy initialization tractable. InterleaveThinker trains planner and critic components separately using step-level supervision rather than full expensive interleaved trajectories~\cite{avg260613679}. DiTTo constructs order-aware restoration trajectories offline and trains a policy from them~\cite{avg260530915}. Such decomposition can simplify optimization and diagnosis, but it also introduces interface risks: an abstract planner may propose unrealizable actions, while a transition-level critic may overlook long-horizon interactions. Supervised competence therefore remains most meaningful when the resulting decisions are validated against executable or rendered consequences.

\textbf{Preference optimization} complements supervised imitation when several trajectories are plausible but their relative quality can be judged. The trajectory pairs described in Section~\ref{subsec:trajectory-supervision} can supervise choices such as whether clarification is needed, how much planning is appropriate, or which repair should be preferred. Clarify Before Executing uses selected trajectories for SFT and preference pairs for direct preference optimization~\cite{avg260716352}. GenEvolve and GenAgent likewise use comparisons between candidate trajectories or consecutive visual states in later optimization or distillation stages~\cite{chen2026genevolve,jiang2026genagent}. Preference objectives capture trade-offs that token-level imitation cannot express directly, but remain sensitive to evaluator bias, generator stochasticity, and user preferences that are absent from the comparison data.

The main limitation of supervised and preference-based initialization is coverage. Demonstrations encode only the tools, failures, visual conditions, and stopping conventions encountered during data construction. An agent may reproduce polished trajectories yet fail when a new artifact, tool, or generator changes the transition dynamics. \textbf{Interaction-based optimization} addresses part of this limitation, but reinforcement learning does not simply replace supervised initialization. SFT establishes executable and structured behavior; preference learning introduces relative trade-offs; subsequent interaction-based learning can then optimize decisions whose value becomes visible only after execution.

\subsection{Reinforcement Learning from Multimodal and Executable Feedback}
\label{subsec:multimodal-rl}

Reinforcement learning is particularly relevant to AVG because the quality of an action often becomes visible only after execution. A plausible prompt rewrite may lose identity, restoration may remove noise while damaging texture, and syntactically valid CAD code may still corrupt geometry. Additional revisions can also reduce an already acceptable result. The learning signal therefore depends on the visual or executable transition produced by an action. Such feedback can be delayed, stochastic, heterogeneous, and vulnerable to evaluator exploitation. The central distinction is therefore the source and temporal granularity of the reward rather than the particular policy-optimization acronym.

\begin{table}[t]
  \caption{Reward sources and failure modes in training \AVG{} policies.}
  \label{tab:reward-sources}
  \small
  \setlength{\tabcolsep}{4pt}
  \renewcommand{\arraystretch}{1.08}
  \renewcommand{\tabularxcolumn}[1]{m{#1}}
  \begin{tabularx}{\textwidth}{@{}
    >{\raggedright\arraybackslash\bfseries}m{0.24\textwidth}
    >{\raggedright\arraybackslash}m{0.27\textwidth}
    >{\raggedright\arraybackslash}m{0.16\textwidth}
    >{\raggedright\arraybackslash}X@{}}
    \toprule
    Reward source & \textbf{What it supervises} & \textbf{Granularity} & \textbf{Main failure mode} \\
    \midrule
    Format and tool-call validity & Action syntax; required fields; tool selection & Action or step & Formal compliance without visual usefulness \\
    \midrule
    Executability and deterministic checks & Code execution; rendering; CAD validity; rule satisfaction & Step or trajectory & Runnable but visually or semantically poor artifacts \\
    \midrule
    Constraint and task rewards & Object, relation, preservation, localization, or task success & Step or outcome & Sparse specifications encourage shortcuts \\
    \midrule
    Perceptual and artifact-quality rewards & Realism, fidelity, restoration, aesthetics, and identity & State, transition, or outcome & Evaluator bias; generator stochasticity; metric conflict \\
    \midrule
    Process and comparative rewards & Whether reflection or repair improves a later state & Transition or trajectory pair & Local gains may not improve the final trajectory \\
    \midrule
    Human or learned preference & Open-ended quality, usefulness, style, and interaction trade-offs & Outcome or trajectory pair & Overoptimization; population mismatch \\
    \bottomrule
  \end{tabularx}
\end{table}
\FloatBarrier

Executable feedback is particularly informative when relevant correctness conditions can be checked mechanically. TOOLCAD uses step- and outcome-level signals from a CAD environment before online curriculum reinforcement learning~\cite{gong2026toolcad}. The Manim study combines code-side and visual signals because successful rendering does not guarantee fidelity to the instruction~\cite{avg260418364}. Similar distinctions arise in web, 3D, and diagram generation. Deterministic checks can establish executability or exact constraint satisfaction, whereas perceptual rewards assess properties that become visible only in the rendered artifact. Treating these signals separately prevents formal validity from masking poor visual or semantic quality.

Outcome rewards align optimization with the final task but create long-horizon credit-assignment problems. GenAgent combines final-image pointwise rewards with pairwise process rewards comparing consecutive states~\cite{jiang2026genagent}. CanvasAgent likewise uses outcome and process signals while adapting tool choices from intermediate assets~\cite{avg260705465}. ImageEdit-R1 combines format and semantic-consistency rewards to coordinate specialized editing agents~\cite{avg260308059}. Denser process supervision can identify useful intermediate decisions, but excessive emphasis on process may reward unnecessarily long traces, whereas outcome-only supervision cannot determine which action caused success or damage.

Generation Navigator exposes a specifically visual form of this problem. A terminal-only reward provides coarse credit even when visual quality peaks early, declines after a later revision, or remains unchanged across wasteful turns. Its PRE-GRPO objective separates discovering a high-quality state, retaining that quality, and reaching it efficiently~\cite{avg260517969}. This suggests that AVG objectives benefit from modeling quality dynamics rather than final quality alone. In iterative creation, improvement, regression, recovery, and no-op behavior can carry different implications even when two trajectories end with similar artifacts.

Perceptual rewards are also noisy because the controller does not fully determine the rendered result. Gen-Searcher finds image-only rewards unreliable when a stochastic generator cannot realize a useful reference, motivating complementary text and image rewards~\cite{avg260328767}. TIR-Agent adapts the contribution of heterogeneous image-quality signals~\cite{avg260327742}. Restore-R1 uses multimodal perceptual feedback~\cite{avg251218599}, while PaAgent balances perceptual, objective, and identity-sensitive criteria~\cite{avg260317055}. Reward composition is therefore most informative when its causal scope matches the component being optimized. Pure outcome rewards may penalize a good decision for an unlucky generator sample, while rewards detached from rendering may reinforce plans that are semantically plausible but visually ineffective.

The verifier or reward model can itself become part of the training problem. JarvisEvo jointly optimizes an editor and a human-calibrated evaluator~\cite{avg251123002}. Its ablations show that increasing internal scores can coexist with weaker external performance when evaluator calibration is inadequate. Updating the evaluator may help maintain a useful learning signal as the policy changes, but it also creates co-adaptation and removes evaluator independence. Claims based on such training are therefore stronger when the learned reward is calibrated against evidence outside the coupled training loop, such as held-out metrics, alternative evaluator families, or human judgments.

Finally, visual reinforcement learning is constrained by rollout cost. Repeated generation, rendering, execution, and multi-agent interaction can dominate training time, motivating offline trajectories, simulation, reward-model proxies, and component-wise optimization. InterleaveThinker separates expensive trajectories into planner and critic stages~\cite{avg260613679}, while VisionCreator performs reinforcement learning in a simulated environment~\cite{avg260302681}. These approximations exchange environmental fidelity for scale and can introduce simulator or reward-model mismatch. Their training value therefore depends on whether policies learned from predicted or proxy consequences transfer back to the deployed generator, tools, and environment.

\subsection{Experience Reuse, Skill Abstraction, and Distillation}
\label{subsec:experience-distillation}

Cross-task improvement is not restricted to gradient-based updates. An AVG system may retain episodes, abstract procedures, update routing knowledge, revise a verifier, or distill successful behavior into model parameters. These mechanisms differ in persistence, scope, and evidential strength. Passive records do not change behavior by themselves. Retrieval can support later decisions without changing model parameters, skills can compress recurring procedures across tasks, and parametric distillation can internalize experience but introduce forgetting. Table~\ref{tab:experience-reuse} organizes these forms of persistent cross-task adaptation.

\begin{table}[t]
  \caption{Levels of experience reuse and persistent cross-task adaptation.}
  \label{tab:experience-reuse}
  \small
  \setlength{\tabcolsep}{5pt}
  \renewcommand{\arraystretch}{1.08}
  \renewcommand{\tabularxcolumn}[1]{m{#1}}
  \begin{tabularx}{\textwidth}{@{}
    >{\raggedright\arraybackslash\bfseries}m{0.28\textwidth}
    >{\raggedright\arraybackslash}m{0.33\textwidth}
    >{\raggedright\arraybackslash}X@{}}
    \toprule
    Update level & \textbf{Persistent change} & \textbf{Minimum evidence} \\
    \midrule
    Passive storage & Logs or artifacts are retained & Records are available for later use \\
    \midrule
    Episodic retrieval & Prior trajectories inform similar later tasks & Retrieval ablation on held-out episodes \\
    \midrule
    Skill abstraction & Multiple traces become reusable procedures or workflows & Skill ablation and transfer across tasks or harnesses \\
    \midrule
    Router or policy-memory update & Tool or model selection changes without base-model fine-tuning & Held-out routing gains with fixed tools and budget \\
    \midrule
    Parametric experience distillation & Experience is internalized in policy parameters & Pre/post-transfer; independent evaluation; update ablation \\
    \midrule
    Continual self-improvement & Repeated updates improve future behavior while retaining prior capabilities & Temporal or task-family splits; forgetting tests; rollback \\
    \bottomrule
  \end{tabularx}
\end{table}
\FloatBarrier

Episodic retrieval provides a basic form of persistent cross-task adaptation without changing the base model. SIDiffAgent stores prior diffusion experience while keeping the underlying model training-free~\cite{avg260202051}. Self-Evolving Agentic Image Restoration indexes high-reward experience by degradation type~\cite{avg260628971}. EvoIR-Agent uses a hierarchical experience pool to guide later tool selection and ordering~\cite{avg260522208}. Such systems can reduce repeated search, but improved performance does not by itself establish learning in the parametric sense. Retrieval-disabled controls and held-out future tasks help distinguish genuine behavioral reuse from simply providing additional context or selected demonstrations.

Skill abstraction increases the scope of reuse by converting multiple traces into reusable procedures. COMFYCLAW distills workflow trajectories, execution errors, and verifier feedback into reusable Agent Skills across agents, backbones, and task splits~\cite{avg260701709}. A skill differs from episodic retrieval because it represents a procedure intended to instantiate in new contexts rather than a particular prior episode. This abstraction also broadens the impact of erroneous updates: a flawed skill can affect many later tasks. Evidence for skill-level adaptation is therefore stronger when transfer extends beyond the examples used to construct the skill and when the system preserves provenance and supports versioning, conflict handling, or rollback.

Router and policy-memory updates lie between episodic retrieval and parametric policy learning. OctoT2I uses self-interaction feedback to update the router's external tool knowledge without supervised fine-tuning~\cite{jiang2026octot2i}. This non-parametric update can alter later generator or tool choices while leaving the router and generator parameters unchanged. Its causal effect is clearest when held-out decisions improve under fixed tools, generation budgets, and retrieval costs. Otherwise, apparent gains may arise from access to additional candidates or more expensive resources rather than better routing knowledge.

GenEvolve provides a contrasting example of parametric visual-experience distillation. It compares tool-orchestrated trajectories, uses a privileged teacher to structure their differences, and updates the agent through on-policy self-distillation and reinforcement learning~\cite{chen2026genevolve}. Instead of retaining expensive search traces as context, the pipeline attempts to internalize their useful behavior into a reusable policy. Evidence for such distillation requires more than parameter change: post-update performance should improve independently of the experience-construction process, and comparison with direct retrieval is needed to determine whether useful decision knowledge has actually been internalized.

The survey's L5 criterion requires more than persistent storage or within-episode correction: retained experience must change later decisions and transfer to held-out tasks. Section~\ref{subsec:training-transfer} examines the transfer, forgetting, budget, and governance controls needed to substantiate this claim.

\subsection{Joint Optimization, Transfer, and Evidence Gaps}
\label{subsec:training-transfer}

AVG systems distribute responsibility across interdependent modules. A planner trained against one executor may exploit its specific strengths and weaknesses, then fail when the generator changes. Independently trained tools may also perform well alone but interact poorly in sequence. OPERA jointly optimizes restoration planning and execution through reinforcement learning and agent-guided tool co-training~\cite{avg260522104}. This design enlarges the reachable behavior space but complicates attribution. Evaluation should separate planner-only, executor-only, and joint updates, using frozen-component and component-swap controls where possible.

Editor-evaluator coupling creates another form of co-adaptation. JarvisEvo updates both policies so the evaluator remains informative as the editor changes~\cite{avg251123002}. This may resist exploitation of a static reward model, but it creates a moving objective and removes evaluator independence. The critic may improve calibration, drift toward editor-specific patterns, or leak reward regularities into the policy. Final evidence should therefore come from held-out metrics, alternative evaluator families, or blinded human assessment. Integrated planner-generator-critic models likewise require role-specific ablations and independent evaluation despite lower interface overhead.

\textbf{Cross-task transfer} remains the principal unresolved question. Most studies evaluate with the same tools, generator family, evaluator family, or task distribution used to construct training trajectories. These experiments demonstrate in-distribution optimization, not reusable agency. Stronger protocols should separate unseen task instances, task families, tool combinations or versions, and generators or executors. They should also distinguish syntactic transfer from decision-quality transfer. An agent may call a new API correctly yet choose it at the wrong time or misread its visual failure modes.

\textbf{Controlled deployment conditions} are also required to substantiate training gains. Comparisons between pre- and post-training policies should match rollouts, generations, retrieved examples, tool calls, and revision turns. Otherwise, a larger test-time budget may appear as a training gain. Efficiency claims should compare both quality at matched cost and cost at matched quality. TIR-Agent and Generation Navigator treat redundant actions and trajectory efficiency as learned behavior~\cite{avg260327742,avg260517969}. These inference quantities serve as causal controls, not as an inference taxonomy.

\textbf{Continual training} exposes failures absent from one-shot fine-tuning. New experience may overwrite prior capabilities, faulty verifiers may promote harmful skills, and growing memories may retrieve stale or conflicting procedures. Evidence should include learning curves, retention on prior tasks, update provenance, conflict handling, and rollback to a known version. User or retrieved assets may also become embedded in shared memory, skills, or parameters, increasing privacy and copyright risk. Bounded update authority is therefore part of training design.

\begin{keytakeaway}[Key Takeaways]
Agentic training depends on faithful trajectories and rewards whose causal scope matches the component being trained. SFT provides executable priors, preference learning encodes relative trade-offs, reinforcement learning connects decisions to multimodal consequences, and experience distillation may convert successful search into persistent capability. Claims of self-improvement require independent evaluation, controlled budgets, transfer beyond experience-construction examples, resistance to forgetting, and governance over retained and updated information.
\end{keytakeaway}
